\chapter{数据结构}\label{ux6570ux636eux7ed3ux6784}

\section{STL 必会组件}\label{stl-ux5fc5ux4f1aux7ec4ux4ef6}

\begin{longtable}[]{@{}lll@{}}
\toprule\noalign{}
需求 & 容器/算法 & 复杂度 \\
\midrule\noalign{}
\endhead
\bottomrule\noalign{}
\endlastfoot
动态数组 & \passthrough{\lstinline!vector!} & 末尾均摊 \passthrough{\lstinline!O(1)!}，随机访问 \passthrough{\lstinline!O(1)!} \\
双端队列 & \passthrough{\lstinline!deque!} & 两端 \passthrough{\lstinline!O(1)!} \\
栈/队列 & \passthrough{\lstinline!stack!} / \passthrough{\lstinline!queue!} & \passthrough{\lstinline!O(1)!} \\
优先队列 & \passthrough{\lstinline!priority\_queue!} & 插入/删除顶 \passthrough{\lstinline!O(log n)!} \\
有序集合 & \passthrough{\lstinline!set!} / \passthrough{\lstinline!multiset!} & 增删查 \passthrough{\lstinline!O(log n)!} \\
有序键值 & \passthrough{\lstinline!map!} & \passthrough{\lstinline!O(log n)!} \\
哈希键值 & \passthrough{\lstinline!unordered\_map!} & 期望 \passthrough{\lstinline!O(1)!}，最坏可能退化 \\
去重/计数 & \passthrough{\lstinline!sort!} + \passthrough{\lstinline!unique!} / \passthrough{\lstinline!map!} & \passthrough{\lstinline!O(n log n)!} \\
\end{longtable}

\passthrough{\lstinline!set!} 的 \passthrough{\lstinline!lower\_bound!}、\passthrough{\lstinline!priority\_queue<pair<...>>!} 的字典序比较必须熟练。

\passthrough{\lstinline!multiset::erase(iterator)!} 与 \passthrough{\lstinline!erase(value)!} 的差别也要注意。

\begin{lstlisting}[language={C++}]
vector<int> a = input;
sort(a.begin(), a.end());
a.erase(unique(a.begin(), a.end()), a.end());
priority_queue<pair<int, int>> pq;
for (int i = 0; i < n; ++i) pq.push({value[i], i});
while (!pq.empty()) {
    auto [value, id] = pq.top(); pq.pop();
    // 处理当前最大值；必要时用 version[id] 判断是否过期
}
\end{lstlisting}

\section{并查集（Disjoint Set Union，DSU）}\label{ux5e76ux67e5ux96c6disjoint-set-uniondsu}

维护若干不相交集合，支持合并和查询连通性。路径压缩 + 按大小合并后，单次摊销 \passthrough{\lstinline!O(α(n))!}。

可扩展：维护集合大小、权值和、到父亲的距离/奇偶、撤销栈。带权并查集常用于维护 \passthrough{\lstinline!dist[x]-dist[y]=w!} 或敌友关系。

\begin{lstlisting}[language={C++}]
struct DSU {
    vector<int> p, sz;
    DSU(int n): p(n + 1), sz(n + 1, 1) { iota(p.begin(), p.end(), 0); }
    int find(int x) { return p[x] == x ? x : p[x] = find(p[x]); }
    bool unite(int a, int b) {
        a = find(a); b = find(b); if (a == b) return false;
        if (sz[a] < sz[b]) swap(a, b);
        p[b] = a; sz[a] += sz[b]; return true;
    }
};
\end{lstlisting}

Kruskal、离线加边、连通块统计、最小生成树、带限制的图连通性都优先考虑 DSU。

\section{树状数组（Fenwick Tree）}\label{ux6811ux72b6ux6570ux7ec4fenwick-tree}

支持单点加、前缀和、区间和，\passthrough{\lstinline!O(log n)!}；空间 \passthrough{\lstinline!O(n)!}。通过差分可支持区间加/单点查；两棵树可支持区间加/区间和。下标必须从 1 开始。

\begin{lstlisting}[language={C++}]
struct BIT {
    int n; vector<long long> t;
    BIT(int n): n(n), t(n + 1) {}
    void add(int x, long long v) { for (; x <= n; x += x & -x) t[x] += v; }
    long long sum(int x) const { long long r = 0; for (; x; x -= x & -x) r += t[x]; return r; }
    long long range(int l, int r) const { return sum(r) - sum(l - 1); }
};
\end{lstlisting}

应用：逆序对、动态排名、离线矩形计数、二维偏序（配合 CDQ/排序）、第 k 个值（二进制提升）。

\section{线段树}\label{ux7ebfux6bb5ux6811}

支持任意可合并区间信息。节点存区间答案，\passthrough{\lstinline!push\_up!} 合并左右儿子；区间修改要有懒标记并在访问子节点前 \passthrough{\lstinline!push\_down!}。

\subsection{常见维护}\label{ux5e38ux89c1ux7ef4ux62a4}

\begin{itemize}
\tightlist
\item
  区间和 + 区间加：标记加法，子节点和增加 \passthrough{\lstinline!tag * len!}。
\item
  区间和 + 区间乘/加：标记为仿射变换 \passthrough{\lstinline!x -> mul*x+add!}，下传时 \passthrough{\lstinline!add = add*mul\_parent + add\_parent!}。
\item
  区间最值 + 区间加：维护最大/最小值及懒加。
\item
  区间赋值：赋值标记覆盖旧加法；标记优先级必须明确。
\item
  扫描线：维护每个离散小区间是否被覆盖，保存 \passthrough{\lstinline!cover!} 与 \passthrough{\lstinline!len!}。
\end{itemize}

复杂度：建树 \passthrough{\lstinline!O(n)!}，单次查询/修改 \passthrough{\lstinline!O(log n)!}。只要操作满足结合律，可把节点信息设计成结构体。

\begin{lstlisting}[language={C++}]
struct SegTree {
    int n; vector<long long> sum, lazy;
    SegTree(int n): n(n), sum(4 * n + 4), lazy(4 * n + 4) {}
    void apply(int p, int l, int r, long long v) {
        sum[p] += v * (r - l + 1); lazy[p] += v;
    }
    void push(int p, int l, int r) {
        if (!lazy[p] || l == r) return;
        int m = (l + r) / 2;
        apply(p * 2, l, m, lazy[p]);
        apply(p * 2 + 1, m + 1, r, lazy[p]);
        lazy[p] = 0;
    }
    void add(int p, int l, int r, int ql, int qr, long long v) {
        if (ql <= l && r <= qr) return apply(p, l, r, v);
        push(p, l, r); int m = (l + r) / 2;
        if (ql <= m) add(p * 2, l, m, ql, qr, v);
        if (qr > m) add(p * 2 + 1, m + 1, r, ql, qr, v);
        sum[p] = sum[p * 2] + sum[p * 2 + 1];
    }
    long long query(int p, int l, int r, int ql, int qr) {
        if (ql <= l && r <= qr) return sum[p];
        push(p, l, r); int m = (l + r) / 2; long long ans = 0;
        if (ql <= m) ans += query(p * 2, l, m, ql, qr);
        if (qr > m) ans += query(p * 2 + 1, m + 1, r, ql, qr);
        return ans;
    }
};
\end{lstlisting}

\subsection{线段树上的二分}\label{ux7ebfux6bb5ux6811ux4e0aux7684ux4e8cux5206}

维护区间和/最大值时，可从根向下找``第一个满足条件的位置''，每层只走一个孩子，\passthrough{\lstinline!O(log n)!}，比外层二分 + 查询更快。

\begin{lstlisting}[language={C++}]
int first_at_least(int p, int l, int r, long long x) {
    if (tree[p] < x) return -1;
    if (l == r) return l;
    int m = (l + r) / 2;
    if (tree[p * 2] >= x) return first_at_least(p * 2, l, m, x);
    return first_at_least(p * 2 + 1, m + 1, r, x);
}
\end{lstlisting}

\section{稀疏表（Sparse Table，ST）}\label{ux7a00ux758fux8868sparse-tablest}

静态数组的幂等运算（min/max/gcd/and/or）预处理 \passthrough{\lstinline!O(n log n)!}，查询 \passthrough{\lstinline!O(1)!}：

\passthrough{\lstinline!k=floor(log2(len))!}，答案为 \passthrough{\lstinline!merge(st[k][l], st[k][r-2\^k+1])!}。不可用于 sum 这类非幂等运算（重叠会重复计数）。

\begin{lstlisting}[language={C++}]
int K = __lg(n);
vector<vector<int>> st(K + 1, vector<int>(n));
st[0] = a;
for (int j = 1; j <= K; ++j)
    for (int i = 0; i + (1 << j) <= n; ++i)
        st[j][i] = min(st[j - 1][i], st[j - 1][i + (1 << (j - 1))]);
auto range_min = [&](int l, int r) {
    int k = __lg(r - l + 1);
    return min(st[k][l], st[k][r - (1 << k) + 1]);
};
\end{lstlisting}

\section{堆与可并堆思想}\label{ux5806ux4e0eux53efux5e76ux5806ux601dux60f3}

\passthrough{\lstinline!priority\_queue!} 用于每步取最小/最大。双堆维护中位数（小根维护大半边，大根维护小半边）。懒删除：堆中保留旧值，弹出时与当前版本/计数比对。

左偏树/可并堆支持两个堆合并，合并复杂度 \passthrough{\lstinline!O(log n)!}；题目要求``合并集合并取最值''时考虑。斜堆实现更短但无严格最坏保证。

\begin{lstlisting}[language={C++}]
priority_queue<int> low; // 小的一半，堆顶是最大值
priority_queue<int, vector<int>, greater<int>> high;
void add_number(int x) {
    if (low.empty() || x <= low.top()) low.push(x); else high.push(x);
    if (low.size() > high.size() + 1) high.push(low.top()), low.pop();
    if (high.size() > low.size()) low.push(high.top()), high.pop();
}
int median() { return low.top(); }
\end{lstlisting}

\section{单调结构与差分约束}\label{ux5355ux8c03ux7ed3ux6784ux4e0eux5deeux5206ux7ea6ux675f}

单调栈/队列见 \href{01-基础与实现规范.md}{01-基础与实现规范}。双端队列还能实现 0-1 BFS：边权 0 从队首加入，边权 1 从队尾加入。

\begin{lstlisting}[language={C++}]
vector<int> dist(n, INF);
deque<int> q;
dist[s] = 0; q.push_front(s);
while (!q.empty()) {
    int u = q.front(); q.pop_front();
    for (auto [v, w] : g[u]) {
        if (dist[v] <= dist[u] + w) continue;
        dist[v] = dist[u] + w;
        if (w == 0) q.push_front(v);
        else q.push_back(v);
    }
}
\end{lstlisting}

\section{字典树（Trie，Prefix Tree）}\label{ux5b57ux5178ux6811trieprefix-tree}

字符集固定时数组儿子最快；动态字符集用 \passthrough{\lstinline!map!}。插入/查询长度 \passthrough{\lstinline!L!} 的字符串均为 \passthrough{\lstinline!O(L)!}。二进制 Trie 支持最大异或、区间内最大异或（结合持久化或离线前缀）。

\begin{lstlisting}[language={C++}]
struct BinaryTrie {
    int ch[200000][2]{}, cnt[200000]{}, nodes = 1;
    void insert(int x) {
        int u = 1; ++cnt[u];
        for (int b = 30; b >= 0; --b) {
            int bit = (x >> b) & 1;
            if (!ch[u][bit]) ch[u][bit] = ++nodes;
            u = ch[u][bit]; ++cnt[u];
        }
    }
    int max_xor(int x) {
        int u = 1, ans = 0;
        for (int b = 30; b >= 0; --b) {
            int bit = (x >> b) & 1;
            if (ch[u][bit ^ 1] && cnt[ch[u][bit ^ 1]])
                ans |= 1 << b, u = ch[u][bit ^ 1];
            else u = ch[u][bit];
        }
        return ans;
    }
};
\end{lstlisting}

\section{可持久化数据结构}\label{ux53efux6301ux4e45ux5316ux6570ux636eux7ed3ux6784}

每次修改只复制从根到叶子的 \passthrough{\lstinline!O(log n)!} 个节点，旧版本仍可查询。主席树（可持久化权值线段树）用于静态区间第 k 小：\passthrough{\lstinline!root[r] - root[l-1]!} 在树上同步下行，\passthrough{\lstinline!O(log V)!}。注意节点池容量约为 \passthrough{\lstinline!版本数 × log V!}。

\begin{lstlisting}[language={C++}]
struct Node { int l = 0, r = 0, sum = 0; } tr[4000000];
int tot = 0;
int update(int old, int l, int r, int pos) {
    int now = ++tot; tr[now] = tr[old]; ++tr[now].sum;
    if (l != r) {
        int m = (l + r) / 2;
        if (pos <= m) tr[now].l = update(tr[old].l, l, m, pos);
        else tr[now].r = update(tr[old].r, m + 1, r, pos);
    }
    return now;
}
int kth(int u, int v, int l, int r, int k) {
    if (l == r) return l;
    int left_count = tr[tr[v].l].sum - tr[tr[u].l].sum;
    int m = (l + r) / 2;
    if (k <= left_count) return kth(tr[u].l, tr[v].l, l, m, k);
    return kth(tr[u].r, tr[v].r, m + 1, r, k - left_count);
}
\end{lstlisting}

\section{分块与莫队算法（Mo's Algorithm）}\label{ux5206ux5757ux4e0eux83abux961fux7b97ux6cd5mos-algorithm}

\subsection{块状数组}\label{ux5757ux72b6ux6570ux7ec4}

把数组分成 \passthrough{\lstinline!√n!} 块，块内暴力、块间整块处理；区间加/区间和约 \passthrough{\lstinline!O(√n)!}。适合操作复杂、难以线段树维护的信息。

\begin{lstlisting}[language={C++}]
int B = max(1, (int)sqrt(n));
vector<long long> a(n), tag((n + B - 1) / B);
void add(int l, int r, long long v) {
    int bl = l / B, br = r / B;
    if (bl == br) { for (int i = l; i <= r; ++i) a[i] += v; return; }
    for (int i = l; i < (bl + 1) * B; ++i) a[i] += v;
    for (int b = bl + 1; b < br; ++b) tag[b] += v;
    for (int i = br * B; i <= r; ++i) a[i] += v;
}
\end{lstlisting}

\subsection{莫队}\label{ux83abux961f}

离线排序询问的左右端点，维护当前区间并增删元素；移动总复杂度约 \passthrough{\lstinline!O((n+q)√n)!}。需要可逆的 \passthrough{\lstinline!add/remove!}，如频次、不同数个数、答案贡献。带修改莫队多一个时间维，块大小约 \passthrough{\lstinline!n\^(2/3)!}。

\begin{lstlisting}[language={C++}]
int block;
struct Query { int l, r, id; };
sort(qs.begin(), qs.end(), [&](const Query& x, const Query& y) {
    int bx = x.l / block, by = y.l / block;
    if (bx != by) return bx < by;
    return (bx & 1) ? x.r > y.r : x.r < y.r;
});
int L = 0, R = -1, distinct = 0;
auto add = [&](int p) { if (++cnt[a[p]] == 1) ++distinct; };
auto remove = [&](int p) { if (--cnt[a[p]] == 0) --distinct; };
for (auto qu : qs) {
    while (L > qu.l) add(--L); while (R < qu.r) add(++R);
    while (L < qu.l) remove(L++); while (R > qu.r) remove(R--);
    answer[qu.id] = distinct;
}
\end{lstlisting}

\section{平衡树思想（了解 Treap）}\label{ux5e73ux8861ux6811ux601dux60f3ux4e86ux89e3-treap}

随机优先级的 BST，支持按键插入、删除、排名、第 k 小、分裂/合并，期望 \passthrough{\lstinline!O(log n)!}。无序序列维护时使用 FHQ Treap，节点存子树大小、区间反转标记、区间和/最大值。需要严格稳定最坏复杂度时优先线段树或 \passthrough{\lstinline!std::set!}。

\begin{lstlisting}[language={C++}]
struct Node { int l, r, key, pri, sz; } tr[MAXN];
int root, tot;
int size(int u) { return u ? tr[u].sz : 0; }
void pull(int u) { tr[u].sz = size(tr[u].l) + size(tr[u].r) + 1; }
void split(int u, int key, int &x, int &y) {
    if (!u) return x = y = 0, void();
    if (tr[u].key <= key) split(tr[u].r, key, tr[u].r, y), x = u;
    else split(tr[u].l, key, x, tr[u].l), y = u;
    pull(u);
}
int merge(int x, int y) {
    if (!x || !y) return x ? x : y;
    if (tr[x].pri > tr[y].pri) return tr[x].r = merge(tr[x].r, y), pull(x), x;
    return tr[y].l = merge(x, tr[y].l), pull(y), y;
}
\end{lstlisting}

\section{撤销与可回滚结构}\label{ux64a4ux9500ux4e0eux53efux56deux6edaux7ed3ux6784}

不做路径压缩的 DSU 可用栈记录每次修改，支持回滚到快照 \passthrough{\lstinline!O(log n)!} 合并、\passthrough{\lstinline!O(1)!} 回滚；配合线段树分治时间解决``边在一段时间内存在''的动态连通性。

\begin{lstlisting}[language={C++}]
struct RollbackDSU {
    vector<int> p, sz; vector<pair<int, int>> hist;
    RollbackDSU(int n): p(n + 1), sz(n + 1, 1) { iota(p.begin(), p.end(), 0); }
    int find(int x) { while (x != p[x]) x = p[x]; return x; }
    int snapshot() { return hist.size(); }
    bool unite(int a, int b) {
        a = find(a); b = find(b); if (a == b) return false;
        if (sz[a] < sz[b]) swap(a, b);
        hist.push_back({b, sz[a]}); p[b] = a; sz[a] += sz[b]; return true;
    }
    void rollback(int snap) {
        while ((int)hist.size() > snap) {
            auto [b, old] = hist.back(); hist.pop_back();
            sz[p[b]] = old; p[b] = b;
        }
    }
};
\end{lstlisting}

\section{李超线段树（Li Chao Segment Tree）}\label{ux674eux8d85ux7ebfux6bb5ux6811li-chao-segment-tree}

维护直线集合在离散/整数横坐标上的最值。每个节点保存当前区间中在中点最优的直线，递归把另一条直线送入可能胜出的子区间；插入和查询均为 \passthrough{\lstinline!O(log V)!}。

\begin{lstlisting}[language={C++}]
struct LiChao {
    using ll = long long;
    struct Line {
        // 这份代码维护“最小值”；空直线用极大截距表示。
        // 若改成最大值，把空值和所有比较方向一起反过来。
        ll k = 0, b = (ll)4e18;
        ll get(ll x) const { return k * x + b; }
    };
    int L, R;             // 允许查询的整数横坐标闭区间 [L,R]。
    vector<Line> tr;      // tr[p]：节点区间内当前保留下来的候选直线。
    LiChao(int L, int R): L(L), R(R), tr(4 * (R - L + 1)) {}

    void add_line(Line nw, int p, int l, int r) {
        // 先比较新线和当前线在端点/中点的优劣。
        int m = (l + r) / 2;
        bool lef = nw.get(l) < tr[p].get(l);
        bool mid = nw.get(m) < tr[p].get(m);

        // 中点更优的线留在当前节点，旧线递归到可能赢的半边。
        if (mid) swap(nw, tr[p]);
        if (l == r) return;

        // 两线在左端和中点的相对优劣发生变化，交点在左半。
        if (lef != mid) add_line(nw, p * 2, l, m);
        else add_line(nw, p * 2 + 1, m + 1, r);
    }
    // 对外接口：只需要传入斜率 k 和截距 b，不要手动传根节点参数。
    void add_line(ll k, ll b) { add_line({k, b}, 1, L, R); }

    ll query(ll x, int p, int l, int r) const {
        // x 所在路径上的每个节点都可能保存一条更优直线，必须全部比较。
        ll ans = tr[p].get(x);
        if (l == r) return ans;
        int m = (l + r) / 2;
        if (x <= m) return min(ans, query(x, p * 2, l, m));
        return min(ans, query(x, p * 2 + 1, m + 1, r));
    }
    // x 必须落在构造时给出的 [L,R] 内；横坐标不连续时先离散化查询点。
    ll query(ll x) const { return query(x, 1, L, R); }
};
\end{lstlisting}

求最大值时把比较符号整体反向；若横坐标不是连续整数，先离散化所有查询点。

\section{区间最值势能线段树（Segment Tree Beats）}\label{ux533aux95f4ux6700ux503cux52bfux80fdux7ebfux6bb5ux6811segment-tree-beats}

``区间取 \passthrough{\lstinline!min!} + 区间和/最大值''不能用普通懒标记，因为只有当前最大值会被压低。维护最大值、次大值、最大值出现次数和区间和；当 \passthrough{\lstinline!x!} 严格大于次大值时可以整段打标记。

\begin{lstlisting}[language={C++}]
struct SegBeats {
    using ll = long long;
    static constexpr ll NEG = -(1LL << 60);
    int n;
    vector<ll> sum;       // 区间和。
    vector<ll> mx;        // 区间最大值。
    vector<ll> second_mx; // 严格小于最大值的次大值。
    vector<int> cnt_mx;   // 最大值出现次数。
    SegBeats(const vector<ll>& a): n((int)a.size() - 1),
        sum(4 * n + 4), mx(4 * n + 4), second_mx(4 * n + 4),
        cnt_mx(4 * n + 4) { build(1, 1, n, a); }

    void pull(int p) {
        // 父节点只保存“最大值这一类”和“次大值这一类”，
        // 不需要保存整个儿子数组，因此 chmin 才能整段快速处理。
        sum[p] = sum[p * 2] + sum[p * 2 + 1];
        if (mx[p * 2] == mx[p * 2 + 1]) {
            mx[p] = mx[p * 2]; cnt_mx[p] = cnt_mx[p * 2] + cnt_mx[p * 2 + 1];
            second_mx[p] = max(second_mx[p * 2], second_mx[p * 2 + 1]);
        } else {
            int x = p * 2, y = p * 2 + 1;
            if (mx[x] < mx[y]) swap(x, y);
            mx[p] = mx[x]; cnt_mx[p] = cnt_mx[x];
            second_mx[p] = max(second_mx[x], mx[y]);
        }
    }
    void build(int p, int l, int r, const vector<ll>& a) {
        if (l == r) {
            // 叶子只有一个最大值，次大值设为负无穷。
            sum[p] = mx[p] = a[l]; second_mx[p] = NEG; cnt_mx[p] = 1;
            return;
        }
        int m = (l + r) / 2;
        build(p * 2, l, m, a); build(p * 2 + 1, m + 1, r, a); pull(p);
    }
    void apply_chmin(int p, ll x) {
        // 调用者保证 second_mx[p] < x < mx[p]，
        // 因而只有 cnt_mx 个最大值会下降，sum 可以 O(1) 更新。
        if (x >= mx[p]) return;
        sum[p] -= (mx[p] - x) * cnt_mx[p];
        mx[p] = x;
    }
    void push(int p) {
        // 父节点的 mx 变小后，儿子的最大值可能仍比父亲大，
        // 把父亲的上界传给儿子即可，不需要递归到底。
        if (mx[p * 2] > mx[p]) apply_chmin(p * 2, mx[p]);
        if (mx[p * 2 + 1] > mx[p]) apply_chmin(p * 2 + 1, mx[p]);
    }
    void chmin(int p, int l, int r, int ql, int qr, ll x) {
        // 无交集或当前最大值已经不超过 x 时，操作没有影响。
        if (qr < l || r < ql || x >= mx[p]) return;
        // x 大于次大值时只有最大值受影响，触发势能线段树的核心优化。
        if (ql <= l && r <= qr && x > second_mx[p]) return apply_chmin(p, x);
        // 否则必须下传并继续递归；删掉这个条件会退化甚至超时。
        push(p); int m = (l + r) / 2;
        chmin(p * 2, l, m, ql, qr, x);
        chmin(p * 2 + 1, m + 1, r, ql, qr, x); pull(p);
    }
    ll query_sum(int p, int l, int r, int ql, int qr) {
        if (qr < l || r < ql) return 0;
        if (ql <= l && r <= qr) return sum[p];
        push(p); int m = (l + r) / 2;
        return query_sum(p * 2, l, m, ql, qr)
             + query_sum(p * 2 + 1, m + 1, r, ql, qr);
    }
    void chmin(int l, int r, ll x) { chmin(1, 1, n, l, r, x); }
    ll query_sum(int l, int r) { return query_sum(1, 1, n, l, r); }
};
\end{lstlisting}

